\section{欧拉函数}

$\varphi(x)$ 表示 $1\sim x$ 中与 $x$ 互质的数的个数。

\subsection{求单一数字 $\varphi(x)$}

\verb|getPhi(x)| 求单一数字的 $\varphi(x)$，复杂度 $\mathcal{O}(\sqrt{x})$：

\begin{lstlisting}
int getPhi(int x){
    int res = 1;
    for (int i = 2;i <= x / i;i++){
        if (x % i == 0){
            res *= i-1;
            x /= i;
            while (x % i == 0){
                res *= i;
                x /= i;
            }
        }
    }
    if (x > 1) res *= x-1;
    return res;
}
\end{lstlisting}

欧拉筛线性求 $\varphi$ 数组，\texttt{phi[i]} 为 $\varphi(i)$，需先调用 \verb|Sieve()|：

\subsection{欧拉筛求 $\varphi(x)$}

\begin{lstlisting}
constexpr int N = 1E7 + 3;
std::bitset<N> np;
std::vector<int> phi,primes;
void Sieve(){
    np[0] = np[1] = 1;
    phi.resize(N,1);
    for (int i = 2;i < N;i++){
        if (!np[i]) {
            phi[i] = i-1;
            primes.push_back(i);
        }
        for (const int& x : primes){
            if (x * i >= N) break;
            np[x * i] = 1;
            if (i % x == 0) {
                phi[i * x] = phi[i] * x;
                break;
            }
            phi[x * i] *= (x - 1) * phi[i];
        }
    }
}
\end{lstlisting}

欧拉函数常用于化简 $\gcd$ 之和。由 $n=\sum_{d\mid n}\varphi(d)$ 取 $n=\gcd(a,b)$ 得
\[ \gcd(a,b)=\sum_{d\mid a,\ d\mid b}\varphi(d). \]

对 $i$ 求和并用 $\sum_{i=1}^{n}[d\mid i]=\left\lfloor\frac{n}{d}\right\rfloor$，得到
\[ \sum_{i=1}^{n}\gcd(i,n)=\sum_{d\mid n}\left\lfloor\frac{n}{d}\right\rfloor\varphi(d). \]

\subsection{欧拉定理（欧拉降幂）}

若 $\gcd(a,m)=1$，则 $a^{\varphi(m)}\equiv 1\pmod{m}$。指数为超大数（字符串给出）时按下式降幂：
\[
a^b \bmod m = \begin{cases}
a^{\,b\bmod\varphi(m)}\bmod m, & \gcd(a,m)=1,\\[4pt]
a^{\,b\bmod\varphi(m)+\varphi(m)}\bmod m, & \gcd(a,m)\ne 1,\ b\ge\varphi(m).
\end{cases}
\]
